
\section{Problem 1}

\paragraph{(a)}

% Given input $G$, the algorithm $\mathcal{A}'$ runs $\mathcal{A}$ on $G$ $t=O(\log n)$ times independently and outputs the majority vote.

Given input $G$, the algorithm $\mathcal{A}'$ runs $\mathcal{A}$ independently $t=O(\log n)$ times on $G$ and outputs the majority vote.
The time complexity of $\mathcal{A}'$  is clearly $O(T(n) \cdot \log n)$. We next analyze the correctness.

% Let $G$ be an arbitrary input graph and let $\chi(G)$ equal 1 of $G$ is a yes-instance (i.e., $G$ contains a $k$-clique), and 0 otherwise.

Let $G$ be an arbitrary input graph, and let $p(G)$ equal $2/3$ if $G$ is a yes-instance (i.e., $G$ contains a $k$-clique), and $1/3$ otherwise. That is, $p(G)$ is the probability that $\mathcal{A}$ output $1$ given $G$.
For every $i \in [t]$ let $x_i$ be the output of $\mathcal{A}$ on $G$ in the $i\th$ invocation. 
Note that $\E[x_i] = p(G)$ for every $i$.  
% Note that for every $i$, it holds that $\E[x_i] = p(G)$.  
% Recall that $A$ errs on $G$ iff $|\frac{1}{t}\cdot x_i - \chi(G)| > 1/3$, which is equivalent to $|\frac{1}{t}\cdot x_i - \chi(G)| > 1/3$.
Additionally, note that $\mathcal{A}'$ errs on $G$ iff $\left| \frac{1}{t}\cdot \sum_{i\in[t]} x_i - p(G) \right| > 1/3$.
By the Chernoff bound, it holds that $t = O(\epsilon^{-2}\cdot \log(1/\delta))$ repetitions suffices to get that $\Pr\left[ \left|\frac{1}{t}\cdot\sum_{i\in[t]} x_i - p(G)\right| > \epsilon \right] \leq \delta$.
Thus, using $\epsilon = 1/3$ and $\delta = 1/n^{10}$, we get that $t = O(\log n)$ repetitions suffices to get that $\mathcal{A}'$ errs on $G$ with probability less than $1/n^{10}$.

% By the checkoff bound, 

% For every graph $G$, let $\xi(G)$ equal 1 of $G$ is a yes-instance (i.e., $G$ contains a $k$-clique), and 0 otherwise.

\paragraph{(b)}

We will show an algorithm $\mathcal{A}''$ such that, if the input $G$ is a yes-instance, then $\mathcal{A}''$ finds a $k$-clique w.p.\ $1/n^9$.
$\mathcal{A}''$ outputs "yes" iff it found a $k$-clique.

$\mathcal{A}''$ partitions the set on nodes $V$ into $k+1$ arbitrary sets of length $n/(k+1)$, denoted $V_1,\ldots V_{k+1}$.
For each $i\in [k+1]$, it runs algorithm $\mathcal{A}'$ on $G$ restricted to $V \setminus V_i$, denoted $G(V \setminus V_i)$.
It then recourses on the first $G(V \setminus V_i)$ that returned yes, or returns "fail" if no call returned yes.
At the base case, if the number of nodes in the graph is $k$, $\mathcal{A}''$ directly checks if the graph is a clique (in $O(k^2) = O(1)$ time), and returns yes iff it is indeed a clique.

\textit{Time complexity.} In each recursive call we make a call to $\mathcal{A}''$ which takes $O(T(n) \cdot \log n)$ time.
Since after $i$ recursion rounds the input shrinks to $(\frac{k}{k+1})^i\cdot n$  
% Since the input shrinks from $n$ to $(k/(k+1))\cdot n$ at each recursion round, 
there are at most $O(\log n)$ recursion rounds.
Thus, the total time complexity is at most $O(T(n) \cdot \log^2 n)$. 

\textit{Correctness.}
Consider any yes-instance $G$. 
Note that if all the call to $\mathcal{A}''$ succeed, then $\mathcal{A}''$ will find a $k$-clique. 
Thus, we only need to show that the probability that one of the calls to $\mathcal{A}''$ fails is at most $1/n^9$.
In the $i\th$ recursion round, we make $k+1$ calls to $\mathcal{A}''$ that each fails w.p.\ at most $\frac{1}{((\frac{k}{k+1})^i\cdot n)^{10}}  $.
By the union bound, the probability that one of the calls fail in one of the rounds is at most
$$
\sum_{i} \frac{k+1}{((\frac{k}{k+1})^i\cdot n)^{10}}  
= \frac{k+1}{n^{10}} \cdot \sum_{i} \frac{1}{(\frac{k}{k+1})^{10\cdot i}}
= \frac{k+1}{n^{10}} \cdot O(1)
\leq \frac{1}{n^9}
$$ 
as desired.

% If $G$ is a no-instance, then w.p. at least $\frac{k+1}{n^{10}} \leq \frac{1}{n^9}$ all the calls to $\mathcal{A}''$ 






% Algorithm $\mathcal{A}''$ will find w.p.\ at least $1/n^9$ the 

% For every $G$, we let 